% GENERATED FILE. Edit canonical lessons, then run npm run generate:books.
% canonical-source-hash: fnv1a64:7fbf487d8ecc6abe
% canonical-lessons: TE-C205-rikardu-ceyu, TE-C205-photo-tiyu, TE-C205-pancuko, TE-C205-kanukko, TE-C205-appu-tisuko

\chapter{To record, To take a photo, To share, To find out, To borrow}
\label{ch:te-to-record-to-take-a-photo-to-share-to-find-out-to-borrow}
\hlchaptermodality{205}

\begin{chapteropening}
\textbf{By the end of this chapter:} \emph{I can say to record, to take a photo, to share, to find out and to borrow.}

Complete the last lesson of chapter 205: I can say to record, to take a photo, to share, to find out and to borrow.
\end{chapteropening}

\section[\texorpdfstring{\te{రికార్డు} \te{చేయు} (rikārḍu cēyu)}{రికార్డు చేయు (rikārḍu cēyu)}]{\te{రికార్డు} \te{చేయు} (rikārḍu cēyu) — to record}

\label{lesson:TE-C205-rikardu-ceyu}



\begin{quote}
Before the new one: say the Telugu for to act, then the Telugu for to draw.
\end{quote}

\subsection*{You'll want to know: \te{రికార్డు} \te{చేయు}}
\textbf{\te{రికార్డు} \te{చేయు}} — \emph{rikārḍu cēyu} — \textquotedblleft{}to record\textquotedblright{}.

\begin{grammarlens}[title={What you've built}]
One more word for this chapter.
\end{grammarlens}

\subsection*{Your turn}
\begin{itemize}
\raggedright
  \item \emph{Say it:} \emph{rikārḍu cēyu}
  \item \emph{Say it:} \emph{rikārḍu cēyu}, once more
  \item \emph{Say it:} say \emph{rikārḍu cēyu}
  \item \emph{Recall:} say the Telugu for to act, then the Telugu for to draw, then say \emph{rikārḍu cēyu} again
\end{itemize}

\begin{tcolorbox}[breakable,colback=teal!4,colframe=teal!35!black,title={Before you move on}]
What does \te{రికార్డు} \te{చేయు} mean? (\textquotedblleft{}To record\textquotedblright{}.) Say it once more.
\end{tcolorbox}
\section[\texorpdfstring{\te{ఫోటో} \te{తీయు} (phōṭō tīyu)}{ఫోటో తీయు (phōṭō tīyu)}]{\te{ఫోటో} \te{తీయు} (phōṭō tīyu) — to take a photo}

\label{lesson:TE-C205-photo-tiyu}



\begin{quote}
Before the new one: say the Telugu for to draw, then the Telugu for to record.
\end{quote}

\subsection*{You'll want to know: \te{ఫోటో} \te{తీయు}}
\textbf{\te{ఫోటో} \te{తీయు}} — \emph{phōṭō tīyu} — \textquotedblleft{}to take a photo\textquotedblright{}.

\begin{grammarlens}[title={What you've built}]
One more word for this chapter.
\end{grammarlens}

\subsection*{Your turn}
\begin{itemize}
\raggedright
  \item \emph{Say it:} \emph{phōṭō tīyu}
  \item \emph{Say it:} \emph{phōṭō tīyu}, once more
  \item \emph{Say it:} say \emph{phōṭō tīyu}
  \item \emph{Recall:} say the Telugu for to draw, then the Telugu for to record, then say \emph{phōṭō tīyu} again
\end{itemize}

\begin{tcolorbox}[breakable,colback=teal!4,colframe=teal!35!black,title={Before you move on}]
What does \te{ఫోటో} \te{తీయు} mean? (\textquotedblleft{}To take a photo\textquotedblright{}.) Say it once more.
\end{tcolorbox}
\section[\texorpdfstring{\te{పంచుకో} (pañcukō)}{పంచుకో (pañcukō)}]{\te{పంచుకో} (pañcukō) — to share}

\label{lesson:TE-C205-pancuko}



\begin{quote}
Before the new one: say the Telugu for to record, then the Telugu for to take a photo.
\end{quote}

\subsection*{You'll want to know: \te{పంచుకో}}
\textbf{\te{పంచుకో}} — \emph{pañcukō} — \textquotedblleft{}to share\textquotedblright{}.

\begin{grammarlens}[title={What you've built}]
One more word for this chapter.
\end{grammarlens}

\subsection*{Your turn}
\begin{itemize}
\raggedright
  \item \emph{Say it:} \emph{pañcukō}
  \item \emph{Say it:} \emph{pañcukō}, once more
  \item \emph{Say it:} say \emph{pañcukō}
  \item \emph{Recall:} say the Telugu for to record, then the Telugu for to take a photo, then say \emph{pañcukō} again
\end{itemize}

\begin{tcolorbox}[breakable,colback=teal!4,colframe=teal!35!black,title={Before you move on}]
What does \te{పంచుకో} mean? (\textquotedblleft{}To share\textquotedblright{}.) Say it once more.
\end{tcolorbox}
\section[\texorpdfstring{\te{కనుక్కో} (kanukkō)}{కనుక్కో (kanukkō)}]{\te{కనుక్కో} (kanukkō) — to find out, to discover}

\label{lesson:TE-C205-kanukko}



\begin{quote}
Before the new one: say the Telugu for to take a photo, then the Telugu for to share.
\end{quote}

\subsection*{You'll want to know: \te{కనుక్కో}}
\textbf{\te{కనుక్కో}} — \emph{kanukkō} — \textquotedblleft{}to find out, to discover\textquotedblright{}.

\begin{grammarlens}[title={What you've built}]
One more word for this chapter.
\end{grammarlens}

\subsection*{Your turn}
\begin{itemize}
\raggedright
  \item \emph{Say it:} \emph{kanukkō}
  \item \emph{Say it:} \emph{kanukkō}, once more
  \item \emph{Say it:} say \emph{kanukkō}
  \item \emph{Recall:} say the Telugu for to take a photo, then the Telugu for to share, then say \emph{kanukkō} again
\end{itemize}

\begin{tcolorbox}[breakable,colback=teal!4,colframe=teal!35!black,title={Before you move on}]
What does \te{కనుక్కో} mean? (\textquotedblleft{}To find out, to discover\textquotedblright{}.) Say it once more.
\end{tcolorbox}
\section[\texorpdfstring{\te{అప్పు} \te{తీసుకో} (appu tīsukō)}{అప్పు తీసుకో (appu tīsukō)}]{\te{అప్పు} \te{తీసుకో} (appu tīsukō) — to borrow (money)}

\label{lesson:TE-C205-appu-tisuko}



\begin{quote}
Before the new one: say the Telugu for to share, then the Telugu for to find out.
\end{quote}

\subsection*{You'll want to know: \te{అప్పు} \te{తీసుకో}}
\textbf{\te{అప్పు} \te{తీసుకో}} — \emph{appu tīsukō} — \textquotedblleft{}to borrow (money)\textquotedblright{}.

\begin{grammarlens}[title={What you've built}]
One more word for this chapter.
\end{grammarlens}

\subsection*{Your turn}
\begin{itemize}
\raggedright
  \item \emph{Say it:} \emph{appu tīsukō}
  \item \emph{Say it:} \emph{appu tīsukō}, once more
  \item \emph{Say it:} say \emph{appu tīsukō}
  \item \emph{Recall:} say the Telugu for to share, then the Telugu for to find out, then say \emph{appu tīsukō} again
\end{itemize}

\begin{tcolorbox}[breakable,colback=teal!4,colframe=teal!35!black,title={Before you move on}]
What does \te{అప్పు} \te{తీసుకో} mean? (\textquotedblleft{}To borrow (money)\textquotedblright{}.) Say it once more.
\end{tcolorbox}
